\bookchapter{Capítulo 13: Recomendações de apps de ferramentas}

\chapterintro{Este capítulo seleciona ferramentas de produtividade, saúde e viagens, e complementa com recursos úteis do sistema e entre dispositivos: limpeza e antispam (como Qingzhuo), opções de bloqueio nativas, alertas de clima e terremotos; acesso remoto ao computador (Sunlogin/ToDesk), escritório com IA (Trae), papel de parede animado, transferência rápida em rede local (LocalSend), área de transferência entre dispositivos, códigos dinâmicos (Aegis/Microsoft Authenticator), agenda do TickTick e digitalização de documentos em papel.}

\booksection{Produtividade: gerencie seu tempo e sua atenção}

Apps de produtividade são o “esqueleto” da vida digital: ajudam a organizar ideias, tarefas e materiais soltos. Veja as recomendações por cenário.

\booksubsection{Lista de tarefas (gerenciamento de tarefas)}

\begin{itemize}
  \item TickTick Público: quem precisa registrar tarefas de trabalho, vida e estudo em um só lugar; um dos destaques desta seção. Diferenciais: sincronização entre celular/computador/tablet/Web; lista + calendário + Pomodoro em um só app; entrada em linguagem natural ("reunião amanhã às 15h", "entregar relatório toda segunda") reconhece horários e recorrências; suporta subtarefas, etiquetas, prioridades, anexos e listas colaborativas. Ideal para gerenciar "agenda + hábitos + revisão diária" de forma integrada.
  \item Como usar melhor o TickTick: | | Caixa de entrada: jogue tudo primeiro em "Hoje/Caixa de entrada" e organize à noite em listas de trabalho, casa e estudo.
  \item Vista de calendário: arraste reuniões e prazos para o calendário e veja conflitos — mais intuitivo do que só a lista de tarefas.
  \item Tarefas recorrentes: remédio, contas e exercícios viram regras automáticas, sem digitar toda vez.
  \item Pomodoro: ative o foco direto na tarefa e veja estatísticas do quanto "realmente concluiu" na semana.
  \item Integração com o celular: diga "lembre-me de..." ao assistente de voz e crie lembretes no sistema; para prazos importantes, use também o calendário nativo como redundância.
\end{itemize}

Todoist Público: usuários minimalistas e multiplataforma. Diferenciais: interface limpa, hierarquia clara de projetos (tarefa subtarefa subtarefa da subtarefa) e pontos de karma para incentivar hábitos.

Microsoft To Do Público: profissionais e usuários de Outlook. Diferenciais: totalmente gratuito e sem anúncios; o "Meu Dia" limpa tudo todo dia e alivia a mente; integração profunda com o ecossistema Microsoft.

Notas/Lembretes nativos Público: usuários leves que não querem instalar apps de terceiros. Diferenciais: experiência nativa, zero atraso, adição por Siri/assistente de voz e integração com widgets do sistema.

\booksubsection{Pomodoro (treino de foco)}

\begin{itemize}
  \item Forest Público: quem não se controla e fica rolando o feed. Diferenciais: ao focar, você planta uma árvore; se sair do app antes da hora, ela morre. Usa "vida virtual" como reforço positivo e permite plantar com amigos.
  \item Tide (Chaoxi) Público: quem precisa de ruído branco para focar/ dormir. Diferenciais: sons naturais (chuva, ondas, floresta) + Pomodoro + meditação em um só app, com interface minimalista e acolhedora.
  \item Princípio: 25 minutos de foco + 5 de pausa. Ajuste para 50/10 ou 90/20 conforme o seu ritmo.
\end{itemize}

\booksubsection{Calendário (visualização do tempo)}

\begin{itemize}
  \item Calendário do sistema Público: todos os usuários. Diferenciais: notificações nativas mais estáveis, consumo mínimo de bateria e adição por voz (Siri/assistente).
  \item Google Calendar Público: usuários com vários aparelhos e quem precisa compartilhar agendas. Diferenciais: a função "Metas" encaixa automaticamente horários de exercício/estudo, e o Gmail reconhece eventos nos e-mails.
\end{itemize}

\booksubsection{Notas (conhecimento organizado)}

\begin{itemize}
  \item Notion Público: quem quer unificar notas, tarefas, bases de dados e wiki; estudantes, freelancers e equipes. Diferenciais: base de conhecimento completa, edição em blocos, páginas pessoais, listas de leitura e quadros de projeto, com rico ecossistema de modelos.
  \item Obsidian Público: gestores de conhecimento, pesquisadores e escritores. Diferenciais: armazenamento local em Markdown, links bidirecionais e grafo de relações — suas notas viram um "segundo cérebro", com dados totalmente seus.
  \item flomo Público: quem quer registrar ideias rápidas sem escrever textões. Diferenciais: notas em cartões, como postar no feed; suporte a etiquetas com \# e linha do tempo para revisitar.
  \item Notas da Apple Público: usuários de iPhone/Mac com necessidades leves. Diferenciais: abre rápido, escaneia documentos, desenha, cria listas e sincroniza pelo iCloud; gratuita e bem integrada ao sistema.
  \item Evernote (Yinxiang) Público: usuários avançados que precisam recortar muitas páginas e arquivar materiais. Diferenciais: extensão de web clipper poderosa, OCR de texto em imagens, marca consolidada e estável, mas o plano premium é pago.
\end{itemize}

\booksubsection{Ferramentas de PDF}

\begin{itemize}
  \item PDF Expert Público: usuários de iPad/iPhone, profissionais que anotam contratos e estudantes. Diferenciais: considerado o melhor editor de PDF no iOS; anotações fluidas, assinatura, unir/dividir e preencher formulários.
  \item WPS Público: quem precisa lidar com Word/Excel/PPT/PDF ao mesmo tempo. Diferenciais: suíte completa, boa compatibilidade com nuvens chinesas e arquivos do WeChat; a versão gratuita já basta para o dia a dia.
  \item Adobe Acrobat Público: quem exige rigor no formato PDF (jurídico, design, publicação). Diferenciais: padrão oficial, conversão para PDF/A impecável, sincronização entre dispositivos e assinatura eletrônica profissional.
\end{itemize}

\booksubsection{Escaneamento (digitalizar documentos em papel)}

Contratos, notas fiscais, trabalhos escolares, anotações de quadro branco, cópias de documentos… a câmera do celular já substitui a maioria dos scanners. Fluxo principal: fotografar bem recortar e corrigir automaticamente exportar PDF/imagem OCR opcional.

\begin{itemize}
  \item CamScanner Público: profissionais e estudantes que escaneiam contratos, recibos e documentos com frequência. Diferenciais: detecção automática de bordas, correção de perspectiva e OCR; gera PDFs multipáginas com sincronização na nuvem.
  \item Microsoft Lens Público: estudantes e quem precisa transformar quadro branco/páginas de livro em Word. Diferenciais: gratuito e sem anúncios, exporta em um toque para OneNote/Word, com boa precisão de reconhecimento.
  \item Escaneamento nas Notas da Apple Público: usuários de iPhone com necessidades leves. Diferenciais: não precisa instalar app; toque no ícone da câmera dentro das Notas, com remoção automática de sombras e assinatura — grátis e prático.
  \item Baimiao / CamScanner Web / "Document Scan" da câmera de marca: a câmera nativa de Huawei, Xiaomi, OPPO e vivo costuma ter modo "Documento/Escaner"; no Android também dá para usar o Google Camera ou o scanner do sistema. Bom para fotografar rapidamente uma página de contrato.
  \item Dicas de fotografia: luz uniforme, celular paralelo ao papel, alinhado aos quatro cantos; documentos multipágina fotografar em ordem e depois juntar em PDF; em documentos sensíveis, atente a marcas d'água e ao local de salvamento, e limpe definitivamente em "Excluídos recentemente" após o uso.
\end{itemize}

\booksection{Saúde e exercícios}

Apps de saúde são o “assistente pessoal” mais subestimado do celular — monitoram continuamente sono, frequência cardíaca e atividade, refletindo o estado de longo prazo melhor do que uma visita ocasional ao médico.

\booksubsection{Exercícios e fitness}

\begin{itemize}
  \item Keep Público: quem treina em casa, iniciantes e quem quer emagrecer/modelar. Diferenciais: muitos cursos gratuitos, de HIIT a ioga; planos de treino gerados por IA; comunidade ativa.
  \item Nike Training Club (NTC) Público: quem quer orientação profissional e treino sistemático. Diferenciais: oficial da Nike, todos os cursos gratuitos, demonstrações com treinadores profissionais e organização clara por músculo/equipamento.
  \item Apple Fitness Público: usuários de Apple Watch e quem se motiva com os anéis de atividade. Diferenciais: três anéis diários (Movimento/Exercício/Em pé), metas visuais e integração perfeita com os dados de frequência cardíaca do relógio.
\end{itemize}

\booksubsection{Monitoramento do sono}

\begin{itemize}
  \item AutoSleep Público: usuários de Apple Watch que querem saber a proporção de sono profundo/leve. Diferenciais: rastreia automaticamente ao dormir com o relógio, sem iniciar/parar manualmente; oferece nota de qualidade e dicas de deitar.
  \item Sleep Cycle Público: quem dorme sem relógio e quer monitorar pelo celular. Diferenciais: analisa viradas e roncos pelo microfone, acorda você na fase de sono leve e torna o despertar menos doloroso.
  \item Mi Fitness Público: usuários de Mi Band/relógio Xiaomi. Diferenciais: gratuito e bem integrado à Mi Band; frequência cardíaca, oxigênio e sono em um só lugar, com ótimo custo-benefício.
\end{itemize}

\booksubsection{Meditação e relaxamento}

\begin{itemize}
  \item Calm Público: profissionais estressados e quem tem dificuldade para dormir. Diferenciais: meditação guiada "Daily Calm" + histórias para dormir (lidas por celebridades); interface bonita e muito conteúdo em inglês.
  \item Tide (Chaoxi) Público: usuários de língua chinesa que precisam de ruído branco. Diferenciais: produto nacional de qualidade; sons naturais + cursos de meditação + Pomodoro, com guiagem em chinês mais próxima.
  \item Little Sleep (Xiao Shuimian) Público: quem precisa de sons para dormir e quer gravar sonhos/roncos. Diferenciais: muitos ruídos brancos combináveis (chuva + ventilador + fogueira), timer para desligar, e gravação de sonhos divertida e útil.
\end{itemize}

\booksubsection{Saúde feminina}

\begin{itemize}
  \item Flo Público: quem registra ciclo menstrual, busca gravidez ou acompanha a saúde da mulher. Diferenciais: previsão precisa de menstruação/ovulação por IA; comunidade com perguntas anônimas; interface refinada.
  \item Meiyou Público: usuárias chinesas que querem conteúdo e comunidade em chinês. Diferenciais: conteúdo localizado rico, com conhecimento adequado a cada fase (ciclo/gravidez/gestação/criação dos filhos).
\end{itemize}

\booksubsection{Lembretes de medicamentos}

\begin{itemize}
  \item Assistente de Medicamentos — Público-alvo: pacientes com doenças crônicas que precisam tomar remédios por muito tempo e filhos que cuidam da medicação de pais idosos. Principal destaque: escaneie o código de barras para identificar o remédio, monte automaticamente a lista de medicamentos, receba alertas na hora certa em forma de pop-up e compartilhe o registro de uso com a família.
\end{itemize}

\booksubsection{Central de dados}

\begin{itemize}
  \item Apple Health / Google Fit Público: todos os usuários de smartphone. Diferenciais: reúne dados de celular, smartband e apps de saúde (passos, frequência cardíaca, sono, ciclo etc.) em um "prontuário de saúde" unificado.
\end{itemize}

\booksection{Viagens e deslocamentos}

Fora de casa, alguns apps bem escolhidos evitam muita dor de cabeça. Veja as recomendações por etapa da viagem.

\booksubsection{Voos e trens}

\begin{itemize}
  \item Umetour (Hanglv Zongheng) — Público-alvo: profissionais que voam com frequência. Principal destaque: dados oficiais da aviação civil, as informações mais confiáveis sobre status de voos, trocas de portão de embarque e avisos de atraso; permite consultar o modelo da aeronave e o mapa de assentos.
  \item 12306 Público: todos que viajam de trem na China. Diferenciais: canal oficial e confiável de compra; suporta lista de espera (compra automática quando alguém devolve); prática para bilhetes estudantis e em grupo.
  \item VariFlight (Feichangzhun) — Público-alvo: quem se importa com a pontualidade dos voos e precisa buscar ou deixar alguém no aeroporto. Principal destaque: previsão de atrasos de voos bem precisa; dá para compartilhar a viagem com a família, que acompanha em tempo real sua localização e o horário de chegada.
\end{itemize}

\booksubsection{Mapas e navegação}

\begin{itemize}
  \item Amap (Gaode) Público: usuários em cidades chinesas, motoristas e quem usa apps de transporte. Diferenciais: trânsito em tempo real preciso, boa capacidade de desviar de engarrafamentos, táxi/carona em um só app e muitas opções de transporte público.
  \item Baidu Maps Público: quem precisa de informações de locais e navegação indoor. Diferenciais: street view / navegação AR, mapas internos detalhados de shoppings e atrações, e muitos pacotes de voz.
  \item Google Maps Público: viajantes internacionais. Diferenciais: cobertura global ampla, avaliações locais confiáveis; street view e mapas offline salvam a viagem no exterior.
\end{itemize}

\booksubsection{Tradução}

\begin{itemize}
  \item Google Tradutor Público: viajantes e estudantes de idiomas. Diferenciais: "Conversa" com tradução bidirecional em tempo real e "Câmera" com tradução instantânea de texto — indispensáveis na viagem; suporta pacotes de idiomas offline.
\end{itemize}

\booksubsection{Conversão de moedas}

\begin{itemize}
  \item Conversor de Câmbio Minimalista (Jijian Huilü) — Público-alvo: quem viaja para o exterior, faz compras internacionais ou compras transfronteiriças. Principal destaque: interface extremamente limpa, mais de 160 moedas, câmbio atualizado em tempo real, comparação de várias moedas ao mesmo tempo e conversão rápida segurando o dedo sobre o número.
\end{itemize}

\booksubsection{Clima}

\begin{itemize}
  \item Clima do sistema / China Weather Network Público: todos, como primeira opção para previsões e alertas oficiais. Diferenciais: apps de clima dos fabricantes e canais oficiais como "China Weather" recebem alertas do serviço meteorológico; chuvas fortes, tufões e calor extremo chegam em banners/notificações, com alta autoridade.
  \item Caiyun Weather Público: quem precisa de previsão de chuva minuto a minuto. Diferenciais: alta precisão em "vai chover nas próximas duas horas"; mapa dinâmico de chuva intuitivo e alertas ativos antes de começar a chover.
  \item Carrot Weather Público: entusiastas de clima que gostam de diversão e muitos widgets. Diferenciais: locução de IA bem-humorada, inúmeros widgets personalizáveis, radar profissional; a versão paga é mais completa.
  \item Pontos de alerta: apps de clima precisam de notificações liberadas para enviarem alertas; os alertas têm quatro níveis — azul, amarelo, laranja e vermelho. Vermelho indica risco meteorológico extremo e você deve ajustar os planos conforme o aviso. Mais formas de ativar na seção 7 deste capítulo «Alertas climáticos / Alertas de desastres naturais / Alertas de terremoto».
\end{itemize}

\booksubsection{Guias de viagem}

\begin{itemize}
  \item Mafengwo Público: quem gosta de roteiros detalhados, viagens de carro ou turismo profundo. Diferenciais: relatos originais dos usuários ricos em conteúdo; roteiros, rotas e dicas para evitar armadilhas; guias disponíveis offline.
  \item Xiaohongshu (RED) Público: jovens que querem experiências reais e lugares/comidas fora do óbvio. Diferenciais: formato visual com fotos, textos e vídeos curtos; conteúdos como "roteiro de 3 dias na cidade X" se atualizam rápido e a busca é boa.
\end{itemize}

\booksection{Introdução a automação}

Automação é uma alavanca: “configura uma vez e economiza sempre”. Pode reduzir a zero as dezenas de operações repetidas do seu dia. Aqui é só uma introdução; práticas e tutoriais detalhados estão no Capítulo 23: Pensamento computacional e introdução à automação.

\begin{itemize}
  \item Atalhos do iOS: ferramenta oficial de automação da Apple, sem programar. Em um toque, executa combos como "ao chegar em casa: ligar Wi-Fi + baixar o brilho + anunciar o tempo".
  \item Automação no Android: cenários nativos (como "Tarefas automáticas" da Xiaomi e "Cenários inteligentes" da Huawei) são os mais fáceis; o Tasker é poderoso, mas tem curva de aprendizado íngreme — para entusiastas.
  \item Ideia central: primeiro identifique "o que você repete todo dia", depois pense "dá para substituir por um gatilho + várias ações?". Exemplos detalhados e recursos avançados no Capítulo 23.
\end{itemize}

\booksection{Navegadores e mecanismos de busca: sua porta de entrada para a informação}

O navegador é a sua janela para a internet; o mecanismo de busca é a bússola para navegar no oceano de informações. Usá-los bem evita muitos desvios.

\booksubsection{Como o navegador funciona:}

Ao digitar um endereço, o navegador: (1)consulta o servidor DNS para descobrir o IP do domínio; (2)envia uma solicitação HTTP/HTTPS ao servidor nesse IP; (3)o servidor devolve o código da página (HTML, CSS, JavaScript); (4)o navegador interpreta o código e renderiza a página que você vê.

\booksubsection{Navegadores recomendados:}

\begin{itemize}
  \item Microsoft Edge: baseado no núcleo Chromium, rápido e compatível. Traz o assistente de IA Copilot, que resume páginas, traduz e conversa. Sincroniza entre dispositivos — recomendo muito como navegador principal.
  \item Safari: padrão nos dispositivos Apple, bem integrado ao sistema e econômico. Boa proteção de privacidade, mas com menos extensões.
  \item Chrome: o mais usado no mundo, com muitas extensões. Porém, consome bastante memória e, no celular, gasta mais bateria.
\end{itemize}

\booksubsection{Escolha do mecanismo de busca:}

\begin{itemize}
  \item Bing: da Microsoft, bem integrado ao Edge. O resumo de IA responde direto, sem só listar links — muito mais eficiente. A busca de imagens também é boa.
  \item Yandex: o maior buscador russo, com busca e busca reversa de imagens muito fortes. Ótimo para achar imagens parecidas e identificar objetos.
  \item Baidu: principal buscador de informações na China, mas com muitos anúncios — exige senso crítico.
  \item Google: o buscador mais preciso do mundo, mas na China exige rede especial para ser usado.
\end{itemize}

\booksubsection{Guia de busca eficiente:}

\begin{itemize}
  \item Use a busca com IA: o resumo do Bing e o Copilot do Edge respondem direto e destilam o essencial de um mar de resultados — muito mais rápido do que clicar link por link.
  \item Identifique desinformação e anúncios: itens marcados como "Anúncio" ou "Patrocinado" são promoções pagas, não resultados orgânicos. Verifique a autoridade da fonte e prefira sites oficiais e veículos conhecidos.
  \item Baixar apps corretamente: no buscador, digite "nome do app + site oficial" (por exemplo, "WeChat site oficial") e prefira resultados com domínio oficial. Não instale software de sites de download de terceiros — costumam vir com programas indesejados.
  \item Como reconhecer o site oficial: para baixar o Windows 11, acesse microsoft.com; para o Chrome, google.com. Conferir se o domínio corresponde à marca é o método mais básico.
\end{itemize}

\booksubsection{Conhecimento sobre domínios:}

O domínio é o “endereço” de um site. Saber lê-lo ajuda a julgar a natureza e a origem da página:

\begin{itemize}
  \item Significado dos sufixos comuns: | | .com — organizações comerciais, o mais comum
  \item .cn — domínio nacional da China
  \item .org — organizações sem fins lucrativos
  \item .net — serviços de rede
  \item .edu — instituições de ensino
  \item .gov — órgãos governamentais
  \item .xyz/.tech/.info — domínios genéricos novos, baratos e de fácil registro: tenha cuidado
\end{itemize}

Domínios por país/região: | | .jp — Japão

.kr — Coreia

.uk — Reino Unido

.de — Alemanha

.tw — Taiwan

.hk — Hong Kong

Dica de segurança: sites .gov e .edu costumam ser mais confiáveis, pois o registro é mais rigoroso. Domínios baratos como .xyz e .top são muito usados em phishing — navegue com atenção.

\booksubsection{Bloqueio de anúncios:}

\begin{itemize}
  \item AdGuard Público: todos os usuários incomodados com pop-ups, anúncios em vídeo e anúncios dentro de apps. Diferenciais: bloqueio em nível de sistema (não só no navegador), filtra banners e rastreamento de splash, economiza dados e acelera o carregamento; a versão paga cobre todos os dispositivos.
\end{itemize}

\booksubsection{Navegadores com foco em privacidade:}

\begin{itemize}
  \item Brave Público: quem se importa com privacidade e não quer ser rastreado por anúncios. Diferenciais: baseado em Chromium, bloqueia anúncios e rastreadores desde o início, traz modo anônimo Tor, é mais rápido que o Chrome e usa menos memória.
  \item DuckDuckGo Público: quem quer "pesquisar uma palavra sem ser perseguido por anúncios". Diferenciais: traz o buscador de mesmo nome (sem histórico), limpa todo o rastro de navegação com um toque ("Fire Button") e tem experiência consistente entre plataformas.
\end{itemize}

\booksubsection{Comparativo de navegadores:}

\begin{center}
\begin{tabularx}{\linewidth}{|X|X|X|}
\hline
\textbf{Navegador} & \textbf{Diferenciais} & \textbf{Público} \tabularnewline
\hline
Microsoft Edge & Núcleo Chromium + Copilot AI integrado, resume páginas / traduz / conversa & Profissionais focados em eficiência e usuários pesados de IA (recomendado como principal) \tabularnewline
Safari & Integração profunda com o ecossistema Apple, econômico, forte em privacidade & Usuários de iPhone/Mac que valorizam bateria e integração \tabularnewline
Chrome & O mais usado no mundo, com o maior ecossistema de extensões & Usuários avançados que precisam de muitas extensões e sincronização entre dispositivos \tabularnewline
Brave & Bloqueia anúncios e rastreadores desde o início, mais rápido & Quem valoriza privacidade e quer bloquear anúncios sem instalar extensões \tabularnewline
DuckDuckGo & Buscador homônimo sem histórico; limpa rastros com um toque & Usuários sensíveis a privacidade que não querem rastreamento personalizado \tabularnewline
\hline
\end{tabularx}
\end{center}

\begin{tcbpractice}{Experimente criar uma automação "relaxar antes de dormir"}
\begin{enumerate}
  \item Crie o gatilho: horário específico (por exemplo, 22:30).
  \item Adicione as ações: baixar o volume > ativar Não Perturbe > tocar playlist de relaxamento > reduzir o brilho.
  \item Salve e teste.
\end{enumerate}

\end{tcbpractice}

\booksection{Ferramentas úteis do sistema: limpeza, antispam e proteção}

Além dos apps de produtividade e saúde, o que mais impacta o dia a dia costuma ser “ser incomodado menos, ter espaço de sobra e não cair em golpe”. Esta seção primeiro lista apps úteis comuns e depois detalha o bloqueio de spam nativo do celular — muita gente usa o aparelho há anos e nunca abriu essas opções.

\booksubsection{Limpeza do sistema e otimização de armazenamento}

\begin{itemize}
  \item Qingzhuo Público: usuários de Android que querem limpar cache, resíduos e pastas vazias de forma leve, sem ser incomodados por "suite de gerenciamento". Diferenciais: tamanho pequeno, sem anúncios; varredura profunda de resíduos, localização de arquivos grandes e análise de uso de armazenamento. Mais discreto que os "acelerar em um toque" comerciais, ideal para manutenção periódica por mãos experientes. Atenção: confira os caminhos antes de limpar; não selecione tudo em WeChat, galeria e documentos; para cache de sistema, prefira a ferramenta nativa da marca e use apps de terceiros só como complemento.
  \item Limpeza nativa de armazenamento Público: todos, especialmente idosos e quem não quer complicações. Diferenciais: Android [Configurações] > [Armazenamento] > [Limpar e acelerar]; iOS [Configurações] > [Geral] > [Armazenamento do iPhone]. O sistema conhece melhor o cache dos próprios apps e tem menor risco de apagar algo importante — deve ser a primeira escolha.
  \item Gerenciamento de armazenamento do WeChat/QQ Público: quem conversa muito. Diferenciais: WeChat [Eu] > [Configurações] > [Geral] > [Armazenamento], com limpeza fina por conversa e tipo de arquivo; mais seguro do que ferramentas externas de "limpar WeChat".
  \item Princípio: se dá para usar o que vem do sistema/app, não use terceiros; se dá para selecionar por caminho/tipo, não use "limpar tudo". O objetivo é liberar espaço, não apagar dados importantes junto.
\end{itemize}

\booksubsection{Apps de antispam e segurança}

\begin{itemize}
  \item Centro Nacional Antifraude Público: todos, especialmente idosos, quem trabalha com finanças e quem compra online ou atende números desconhecidos. Diferenciais: oficial; aviso de chamadas, alertas de risco em SMS e detecção de sites suspeitos; integrado às bases antifraude da polícia — a camada mais confiável das "três linhas de defesa".
  \item Centro de segurança da marca / gerenciador do celular Público: usuários da respectiva marca. Diferenciais: Huawei Phone Manager, Mi Security Center, OPPO/vivo Phone Manager etc., com bloqueio de spam, antivírus, gerenciamento de permissões e proteção de pagamentos. Costumam ser gratuitos e integrados ao sistema; evite instalar um segundo pacote comercial para não gerar conflitos de permissão.
  \item Serviço antispam da operadora Público: quem é bombardeado por chamadas de alta frequência e voz por IA. Diferenciais: China Mobile envie KTFSR para 10086; China Telecom siga "Tianyi Antispam"; China Unicom configure no "Unicom Mobile Manager". O bloqueio acontece na operadora e evita até que o celular toque.
  \item AdGuard Público: todos os usuários incomodados com pop-ups, anúncios em vídeo e anúncios dentro de apps. Diferenciais: bloqueio em nível de sistema (não só no navegador), filtra banners e rastreamento de splash, economiza dados e acelera o carregamento; a versão paga cobre todos os dispositivos.
  \item Cuidado com pacotes de "limpeza/antispam": alguns gerenciadores de terceiros ficam rodando em segundo plano e pedem acesso a contatos e SMS sem necessidade. A proteção nem sempre é melhor que a do sistema e ainda gasta bateria e privacidade. Priorize o nativo e, se precisar, some o app oficial antifraude.
\end{itemize}

\textbf{Princípios para escolher apps úteis}\par

\begin{itemize}
  \item Sistema primeiro: bloqueio, limpeza e permissões — ative antes o que já vem no celular.
  \item Complemento oficial: Centro Nacional Antifraude e antispam da operadora são canais oficiais e podem ser somados.
  \item Terceiros com moderação: um app por função já basta; evite conflitos de permissão e disputa de recursos.
  \item Menos é mais: quanto mais ferramentas, mais notificações e permissões — e maior a chance de ser incomodado pelas próprias ferramentas.
\end{itemize}

\booksubsection{Bloqueio de spam nativo: como ativar em cada marca}

A seguir, os caminhos por marca. Os nomes podem variar conforme a versão do sistema; pesquise “spam”, “bloqueio” ou “chamadas desconhecidas” nas configurações. Complementa o Capítulo 8 «Combate a chamadas indesejadas»: aqui focamos nas opções; lá, na marcação de números e recursos por engano.

\begin{center}
\begin{tabularx}{\linewidth}{|X|X|X|}
\hline
\textbf{Plataforma} & \textbf{Caminho} & \textbf{Opções recomendadas} \tabularnewline
\hline
iPhone / iOS & {}[Configurações] > [Telefone] / [Mensagens] & Silenciar chamadas desconhecidas; contatos bloqueados; "Filtrar remetentes desconhecidos" nas Mensagens \tabularnewline
Xiaomi / Redmi & App Telefone > Configurações > Bloqueio de spam; ou Centro de Segurança & Bloqueio inteligente, bloqueio de "toque e desliga", filtro de palavras-chave em SMS \tabularnewline
Huawei / Honor & Telefone > Bloqueio de spam; Gerenciador do celular > Bloqueio de spam & Identificação de chamadas, lista negra, bloqueio inteligente de SMS \tabularnewline
OPPO / realme / OnePlus & App Telefone > Bloqueio de spam; ou Este aparelho > Bloqueio de spam & Bloqueio de chamadas, identificação de golpes, regras de bloqueio de SMS \tabularnewline
vivo / iQOO & Telefone > Bloqueio de spam; i Manager & Bloqueio de spam, alerta de golpes, filtro de mensagens \tabularnewline
Samsung & Telefone > Configurações > Bloquear números / Identificação de chamadas e spam & Bloquear números desconhecidos, identificação de spam (varia por região) \tabularnewline
\hline
\end{tabularx}
\end{center}

\begin{tcbinfo}{Como combinar as três camadas de proteção}
No dia a dia, recomendo “bloqueio do sistema + antispam da operadora + Centro Nacional Antifraude”: o sistema cuida de marcação e lista negra; a operadora bloqueia chamadas de alta frequência na rede; o centro antifraude completa com alertas oficiais. Não é preciso deixar tudo no nível máximo — senão entrega, comida e o médico também são bloqueados. Libere contatos e contatos recentes.

\end{tcbinfo}

\begin{tcbpractice}{Prática: ative o bloqueio de spam hoje}
\begin{enumerate}
  \item Abra as configurações do app de telefone, encontre "Bloqueio de spam / Silenciar chamadas desconhecidas" e ative o bloqueio inteligente.
  \item Confira também o celular dos seus pais: mesmo caminho + antispam da operadora (China Mobile: envie KTFSR para 10086).
  \item Instale o "Centro Nacional Antifraude" na loja de apps, registre-se e ative o aviso de chamadas.
  \item Ligue para um número de teste já marcado (ou veja o registro de bloqueios) para confirmar que funciona.
\end{enumerate}

\end{tcbpractice}

\booksection{Alertas climáticos, de desastres naturais e de terremotos}

O bloqueio de spam protege contra golpes; os alertas protegem contra desastres e clima extremo. Muita gente só instala o app de clima para ver a temperatura e não liga notificações nem alertas de terremoto. Nos segundos ou minutos antes de uma enchente vermelha, um tufão ou um terremoto, o celular pode ganhar tempo precioso para você se proteger.

\booksubsection{Alertas climáticos: chuva forte, tufões e calor extremo}

\begin{itemize}
  \item Priorize canais oficiais: alertas da Administração Meteorológica da China, do Centro Meteorológico Nacional e das estações provinciais são os mais confiáveis. Siga a conta oficial "China Weather" no WeChat e Weibo, ou baixe apps oficiais como "China Weather".
  \item Apps de clima do sistema: Huawei Weather, Mi Weather, OPPO Weather, vivo Weather, Clima da Apple etc. A maioria envia alertas da estação local; abra as configurações do app e confirme que "Alertas meteorológicos / Clima severo" está ativado.
  \item Previsão minuto a minuto (Caiyun etc.): útil para decidir "preciso de guarda-chuva nos próximos 30 minutos" no trajeto ou ao buscar as crianças; a fonte oficial de alertas continua sendo o serviço meteorológico.
  \item Significado das cores dos alertas (geral): | | Azul: atenção, rotina normal, acompanhe atualizações.
  \item Amarelo: redobre a atenção e reduza atividades ao ar livre.
  \item Laranja: grave, fique em casa se possível e afaste-se de áreas de risco.
  \item Vermelho: muito grave, siga a suspensão de aulas/trabalho e as orientações de abrigo. Não atravesse alagamentos, áreas de enxurrada ou trechos com ventania.
\end{itemize}

Tipos comuns de alerta: chuva forte, nevasca, tufão, calor extremo, onda de frio, ventania, raios, granizo, pista congelada, nevoeiro, seca e risco de deslizamentos.

\booksubsection{Alertas de desastres naturais: não é só “chuva”}

\begin{itemize}
  \item Quem publica: a estação meteorológica cuida do clima; deslizamentos e enxurradas costumam ser emitidos em conjunto por recursos naturais, recursos hídricos e defesa civil; governos locais também divulgam por rádio de emergência, SMS e canais como o 12379.
  \item Como receber por vários canais: | | Notificações do sistema + alertas do app de clima.
  \item Siga no WeChat as contas oficiais da "Defesa Civil" e da "Estação Meteorológica" locais.
  \item Canais do sistema nacional de alertas 12379 (SMS/serviço, conforme a região).
  \item TV, rádio de emergência, alto-falantes da vila/comunidade — no campo ou na cheia, não fique só no celular.
\end{itemize}

Cuidados ao ar livre: em tufão, não se aproxime de placas nem passe sob árvores; em chuva forte, evite garagens subterrâneas, túneis e viadutos; no turismo em montanhas, diante de alerta de risco geológico, afaste-se imediatamente de encostas e vales.

Integração com o guia de primeiros socorros: se ficar incomunicável, sem bateria ou precisar pedir socorro com localização, veja o Capítulo 32 (cenas de terremoto/desastre) e o Capítulo 7 (comunicação via satélite e sem rede).

\booksubsection{Alerta de terremoto: não é o mesmo que “previsão de terremoto”}

O alerta de terremoto (EEW) usa o fato de que ondas eletromagnéticas viajam mais rápido que as ondas sísmicas, e que estações captam primeiro a onda P, para avisar a área afetada com alguns a dezenas de segundos de antecedência — antes da onda destrutiva (S) chegar. Ele não prevê se um terremoto vai acontecer; só ganha tempo para se abrigar, cortar o fogo e sair do elevador. Quanto mais perto do epicentro, menor o tempo de aviso — e às vezes não sobra tempo nenhum.

\begin{itemize}
  \item Soluções oficiais e do setor: na China continental, incluem capacidades da rede nacional de alerta sísmico e serviços impulsionados por institutos como o Instituto de Redução de Desastres de Chengdu; alguns fabricantes pré-instalam ou recebem alertas via internet.
  \item Apps de terceiros: procure na loja apps de "alerta de terremoto" de instituições confiáveis (confira o desenvolvedor e evite apps falsos que pedem permissões demais).
  \item O que fazer ao receber o alerta: abrigue-se perto (sob a mesa, canto de parede estrutural, móveis firmes) e proteja cabeça e pescoço; afaste-se de janelas, paredes externas e móveis altos; não corra escadas nem pule; se estiver no elevador, saia no andar mais próximo; ao ar livre, vá para área aberta, longe de fios, placas e construções precárias. Depois, verifique gás e eletricidade e atenção às réplicas.
  \item Não silencie: permita que o alerta de terremoto "fuja" do modo silencioso/Não Perturbe; à noite, deixe o celular ao alcance da mão e não zere o volume.
\end{itemize}

\booksubsection{Alerta de terremoto nativo: como ativar por marca}

Os caminhos podem variar conforme a versão; pesquise “terremoto”, “alerta” ou “alerta de segurança” nas configurações. O iOS tem suporte limitado na China e costuma exigir um app confiável de alerta, com permissões de localização e notificação liberadas.

\begin{center}
\begin{tabularx}{\linewidth}{|X|X|X|}
\hline
\textbf{Plataforma} & \textbf{Caminho comum} & \textbf{Observações} \tabularnewline
\hline
Xiaomi / Redmi & Gerenciador / Centro de Segurança > Alerta de terremoto; ou pesquise "terremoto" nas Configurações & Algumas versões já incluem o serviço; exige internet e localização ativa \tabularnewline
Huawei / Honor & Configurações > Segurança > Notificações de emergência / Alerta de terremoto & O nome pode ser "Notificações de emergência", incluindo terremotos \tabularnewline
OPPO / realme / OnePlus & Configurações > Segurança (SOS) > Alerta de terremoto & Depende do modelo e da cobertura regional do serviço \tabularnewline
vivo / iQOO & App de Clima ou Configurações > Segurança/Mais > Alerta de terremoto & Confira o menu real da sua versão do sistema \tabularnewline
iPhone / iOS & Costuma não haver opção nativa completa; instale um app confiável e permita notificações & O app precisa rodar em segundo plano e ter localização; não bloqueie seus alertas no Não Perturbe \tabularnewline
Samsung etc. & Configurações > Segurança e emergência / Alertas de emergência sem fio & A cobertura dos "alertas de emergência" varia por região e operadora \tabularnewline
\hline
\end{tabularx}
\end{center}

\begin{tcbwarn}{Erros comuns sobre alerta de terremoto}
\begin{itemize}
  \item Alerta $\neq$ previsão: ele não diz "vai ter terremoto amanhã"; só avisa quando o terremoto já aconteceu e as ondas ainda não chegaram.
  \item Nem todo lugar ganha muitos segundos: perto do epicentro podem ser poucos segundos ou nada; o benefício principal é para as áreas vizinhas.
  \item Notificações desligadas = sem alerta: por melhor que seja o app, com som mudo, processos fechados e sem localização, você pode não receber.
  \item Apps falsos: desconfie de "previsão de terremoto" que pede contatos e SMS. Prefira o nativo do sistema ou soluções de instituições sérias.
\end{itemize}

\end{tcbwarn}

\begin{tcbinfo}{Como combinar os alertas}
Sugestão: “alertas do app de clima + canais oficiais de meteorologia/emergência + alerta de terremoto (sistema ou app confiável)”. Como no antispam, não instale uma pilha de apps iguais; o essencial é notificação liberada, localização permitida e o Não Perturbe deixando o alerta passar. Depois de configurar o antispam no celular dos seus pais, confira também os alertas de clima e terremoto.

\end{tcbinfo}

\begin{tcbpractice}{Prática: revise os alertas em 10 minutos}
\begin{enumerate}
  \item Abra o app de clima do sistema, entre nas configurações e ative "Alertas meteorológicos / Clima severo".
  \item Nas configurações, pesquise "terremoto" ou "alerta" e ative o alerta de terremoto/emergência (Android); no iPhone, instale um app confiável e permita notificações.
  \item Siga a conta oficial da estação meteorológica ou defesa civil local e permita as notificações.
  \item Confirme que o som do alerta não fica totalmente mudo no Não Perturbe; faça o mesmo no celular dos seus pais.
\end{enumerate}

\end{tcbpractice}

\booksection{Colaboração entre dispositivos: acesso remoto, transferência e segurança de contas}

O celular flui melhor quando forma equipe com o computador e o tablet. Esta seção traz ferramentas de uso frequente: controlar o PC pelo celular, escritório com IA e colaboração remota, papel de parede animado, transferência rápida em rede local, área de transferência entre dispositivos, códigos 2FA e a ponte entre digitalização e agenda. Para integração mais completa, veja o Capítulo 17.

\booksubsection{Controlar o computador pelo celular: Sunlogin / ToDesk}

\begin{itemize}
  \item Sunlogin Público: quem precisa operar o Windows ou Mac de casa/trabalho de longe; também usado para ajudar idosos remotamente (veja o Capítulo 27). Diferenciais: ecossistema chinês maduro; o app do celular controla a área de trabalho, transfere arquivos e pode ligar/desligar o PC (algumas funções exigem hardware); a versão gratuita costuma bastar para uso pessoal leve.
  \item ToDesk Público: quem quer conexão estável e baixa latência para trabalho/operacoes remotas. Diferenciais: conexão por código de dispositivo + senha, celular controla o PC e vice-versa (conforme a versão), transferência de arquivos e troca de monitores ágeis; boa experiência com nós na China.
  \item Uso básico: instale o cliente no computador e entre/mantenha online instale o app no celular use o código do dispositivo + senha de conexão (ou a mesma conta) opere a área de trabalho com touchpad/gestos no celular.
  \item Limites de segurança (guarde bem): | | Se um desconhecido se passar por "atendente/polícia/operadora" e pedir para instalar Sunlogin ou ToDesk e acessar seu celular ou PC remotamente, recuse — é golpe clássico (veja o Capítulo 27).
  \item Para uso próprio: ative verificação em duas etapas, troque a senha padrão e, quando não usar, saia da conta ou desative o acesso sem supervisão; no computador da empresa, siga as regras de segurança.
  \item Em locais públicos, observe se alguém está espiando a tela ou anotando senhas durante o acesso remoto.
\end{itemize}

\booksubsection{Escritório com IA: Trae / WorkBuddy / Dumate}

\begin{itemize}
  \item Trae Público: quem programa e colabora em documentos, com bom suporte em chinês. Diferenciais: IDE/assistente de programação com IA da ByteDance; o modo SOLO gera projetos a partir de uma frase; muito eficiente para código e documentos. Rodando no computador, dá para acompanhar as tarefas pelo celular com acesso remoto ou os recursos oficiais do produto (consulte a versão atual). Veja o Capítulo 41 (panorama de ferramentas de IA para programação).
  \item WorkBuddy Público: profissionais que querem IA para organizar atas, e-mails, relatórios e documentos. Diferenciais: assistente de IA focado em escritório, ótimo para "redigir — revisar — resumir"; parte dos fluxos pode ser iniciada no celular e o trabalho pesado fica no computador.
  \item Dumate Público: equipes e pessoas que precisam de colaboração, documentos e gestão de conhecimento integrados. Diferenciais: ferramenta de escritório colaborativo com IA, forte em reunir materiais e edição conjunta; no celular dá para ver e editar levemente; diagramação complexa é melhor no PC.
  \item Como entender o "controle remoto do PC pelo celular": muitos recursos de IA rodam no computador ou na web. Com Sunlogin/ToDesk ou o app/site oficial, você "comanda" tarefas de IA no PC durante o trajeto. Princípio: cálculo pesado e documentos longos no computador; acompanhar progresso e ajustar frases, no celular.
  \item Privacidade: antes de enviar contratos, documentos de identidade ou dados de clientes, leia as políticas de privacidade e as exigências de compliance da empresa; para material corporativo, prefira as ferramentas indicadas pela empresa.
\end{itemize}

\booksubsection{Wallpaper animado: N0va Desktop (celular + PC)}

\begin{itemize}
  \item Hoyopaper (Rengong Zhuomian) — Público-alvo: fãs dos personagens da miHoYo que querem papéis de parede animados de graça. Principal destaque: aplicativo gratuito de papéis de parede animados lançado pela miHoYo, disponível para celular e PC; as cenas animadas/Live2D dos personagens podem virar papel de parede da área de trabalho e, em alguns modelos e sistemas, da tela de bloqueio ou da tela inicial.
  \item Bateria e desempenho: wallpaper animado consome mais bateria e GPU que imagem estática. Em bateria baixa ou aparelhos antigos, use só na tomada ou reduza a taxa de quadros (se o app permitir).
  \item Origem da instalação: prefira o site oficial ou lojas confiáveis; evite "versões modificadas" com cavalos de Troia. No PC, confira o caminho de instalação e se o início automático é necessário.
\end{itemize}

\booksubsection{Transferência rápida em rede local: LocalSend / Feige / AirDrop}

\begin{itemize}
  \item LocalSend Público: quem mistura Android/Windows/iOS/Mac/Linux e quer transferir entre plataformas sem passar pela nuvem pública. Diferenciais: código aberto, conexão direta no mesmo Wi-Fi para fotos, vídeos e documentos, sem cabo e sem servidores do fabricante; a velocidade depende do roteador.
  \item Feige Chuanshu (Feige) — Público-alvo: computadores Windows de escritório e cenários antigos de trabalho em rede local. Principal destaque: ferramenta tradicional de transferência de arquivos e de chat em rede local, boa para mandar arquivos grandes pela intranet; atenção: baixe de canais oficiais e siga as políticas de TI na rede da empresa.
  \item AirDrop Público: usuários de iPhone, iPad e Mac. Diferenciais: compartilhamento em um toque no ecossistema Apple, com Bluetooth + Wi-Fi direto e sem gastar dados; configure "Somente contatos" para evitar estranhos.
  \item Outras opções: Mi Share, Huawei Share, vivo Share etc. são muito rápidos entre aparelhos da mesma marca; o "Arquivos" do WeChat é prático, mas tem limite de tamanho e de dispositivos. Para arquivos grandes, prefira LocalSend/AirDrop.
  \item Dicas: ambos no mesmo Wi-Fi; em Wi-Fi público, evite arquivos sensíveis e prefira o hotspot pessoal; ao final, confirme no destinatário se o arquivo está íntegro.
\end{itemize}

\booksubsection{Área de transferência entre dispositivos: KDE Connect / soluções de marca}

\begin{itemize}
  \item Área de transferência universal da Apple Público: usuários de iPhone + Mac. Diferenciais: com o mesmo Apple ID e Handoff ativos, copie no celular e cole no Mac (e vice-versa), para textos e alguns arquivos.
  \item KDE Connect Público: Android + Windows/Linux/Mac, quem quer sincronização de área de transferência open source, espelhamento de notificações e entrada remota. Diferenciais: após parear na LAN, sincroniza a área de transferência, envia notificações do celular ao PC e usa o celular como controle de slides/touchpad; os dados não passam por contas de terceiros.
  \item Soluções de marca: Huawei Multi-Screen Collaboration/Super Device, Mi Share+, vivo Share, OPPO Connect, Honor Connect etc. costumam ter área de transferência entre telas e arrastar arquivos; a experiência é melhor entre aparelhos da mesma marca.
  \item Canal oficial Windows + Android: "Phone Link / Link to Windows" (evolução do Your Phone) mostra notificações e fotos do celular no PC e permite copiar/colar em cenários compatíveis; modelos como Samsung têm integração mais profunda.
  \item Segurança: a área de transferência pode guardar senhas e códigos; em computadores compartilhados, evite a sincronização global e não cole duas vezes em campos de senha. Alguns teclados também leem a área de transferência — atente às permissões.
\end{itemize}

\booksubsection{Códigos dinâmicos 2FA: Aegis / Microsoft Authenticator}

\begin{itemize}
  \item Por que um autenticador em vez de SMS: o código por SMS pode sofrer sequestro de SIM; o TOTP é gerado localmente por algoritmo de tempo, não depende do sinal do número e resiste melhor a phishing e interceptação (ainda assim, proteja-se contra malware e capturas de tela).
  \item Aegis Authenticator Público: usuários de Android que preferem código aberto e controle local. Diferenciais: gratuito e open source, backup criptografado, importação/exportação, interface limpa e desbloqueio da lista de tokens por digital/rosto.
  \item Microsoft Authenticator Público: usuários de contas Microsoft, Office 365/empresariais; iOS/Android. Diferenciais: aprovação de login por push, monitoramento de senhas e backup de contas (conforme a versão), bem integrado ao Windows e à identidade corporativa.
  \item Google Authenticator / 2FAS etc.: também são comuns; o importante é ter backup criptografado para não perder os tokens na troca de aparelho.
  \item Na troca de celular: migre do aparelho antigo ou use os "códigos de backup/chave de recuperação" do serviço. Imprima os códigos em papel e guarde em dois lugares — não só no mesmo celular.
  \item Com gerenciadores de senhas: alguns cofres também guardam TOTP, o que é prático, mas concentra risco. Para contas de alto valor, prefira "gerenciador de senhas + autenticador independente". Veja o Capítulo 14.
\end{itemize}

\booksubsection{Combinação recomendada: como encaixar tudo no dia}

\textbf{Exemplo de fluxo prático}\par

\begin{itemize}
  \item Manhã: veja a agenda no TickTick confira clima/alertas abra Trae/documentos no PC e deixe o celular de lado.
  \item Dia: digitalize documentos em papel com apps de scan e envie via LocalSend/AirDrop para o PC; textos longos, redija no computador com ferramentas de IA.
  \item Fora de casa: conecte ao PC com ToDesk/Sunlogin e continue a tarefa; para logins importantes, use códigos do Aegis/Microsoft Authenticator.
  \item Noite: marque as tarefas concluídas no TickTick; se quiser, troque o wallpaper animado para relaxar.
\end{itemize}

\begin{tcbwarn}{Regras de ouro para acesso remoto e segurança de contas}
\begin{itemize}
  \item Se alguém pedir para instalar software de acesso remoto desligue e procure você mesmo o atendimento oficial.
  \item Nunca informe códigos de verificação ou tokens a ninguém, nem a "entregadores" ou "atendentes".
  \item Antes de controlar o PC da empresa remotamente ou instalar software de IA, pergunte ao TI se é permitido.
  \item Em redes públicas, redobre a atenção ao transferir arquivos ou fazer login remoto.
\end{itemize}

\end{tcbwarn}

\begin{tcbpractice}{Prática: teste a ponte "entre dispositivos" uma vez}
\begin{enumerate}
  \item Conecte celular e PC ao mesmo Wi-Fi e envie uma foto pelo LocalSend (ou AirDrop).
  \item Crie no TickTick uma tarefa "hoje às 20: organizar scans" e confirme que aparece no celular e no PC.
  \item Ative 2FA no e-mail ou WeChat (se disponível), vincule com Aegis ou Microsoft Authenticator e guarde os códigos de backup.
  \item Instale o Sunlogin ou ToDesk no seu computador, conecte uma vez pelo seu celular e depois saia ou reforce a senha — aprenda o fluxo sem deixar brechas.
\end{enumerate}

\end{tcbpractice}
